
\section{The Free Church’s Second Annual Meeting}
June 8–12, 1898

——

\subsection{Wednesday Morning}

The Lutheran Free Church assembled for its annual meeting in Trinity Church, Pastor Gjertsen’s charge, Minneapolis, Minnesota, on June 8, 1898, at 10:30 in the morning.

The meeting was opened with the singing of Hymn 433, verses 1–5. Pastor Chr. Nrehus then delivered the opening sermon on 1 Corinthians 3:10–17.

Now that, after the passing of a year, we are again assembled for an annual meeting here in this church, it is surely with the desire and prayer in our hearts that the meeting may become a blessing to us and to our Church. But if this is to happen, then God himself must be allowed to lead us by his Spirit and by his Word. Only then are we enabled to do the Lord’s work in truth.

According to our text, we will speak a little about the Lord’s temple and its building up.

First, the foundation. When one is to build, it is of great importance that a good foundation be laid; for if one lays a poor foundation, one risks the house falling down and oneself being buried beneath its ruins. So also with the congregation of God. If a good foundation is not laid, the whole structure will collapse.

Now there is only one foundation upon which a true congregation of the Lord can be built, namely Jesus Christ. He is the great chief cornerstone upon which we all, as living stones, must be built up, so that together we may constitute a holy temple for the Lord. It was upon this foundation that the Apostle Paul built. It was the same foundation upon which Luther built when, as the foundation of his life’s work, he laid this truth: that the human being is justified by faith in Christ alone.

Christ, then, is the foundation; but let each one take heed how he builds upon it. It is this admonition of the Apostle that we especially wish to lay upon one another’s hearts at this hour. It is an exceedingly important work, this work of taking part in the building of God’s temple, and we do well to examine ourselves, whether we build well and whether we use the right material.

The Apostle names several things as materials—gold, silver, costly stones, wood, hay, straw. He assumes that all these different materials will be used. He does not here specifically commend any one of them; yet it is nevertheless evident from the context and from other statements of the same Apostle that the right material is that which is here called gold, silver, and costly stones. He says in another of his letters: Come to him, the living Stone, indeed rejected by men, but chosen and precious before God, and yourselves, as living stones, be built up into a spiritual house, a holy priesthood, to offer spiritual sacrifices acceptable to God through Jesus Christ.

The right material, then, is living stones, stones that have been made alive, cleansed, and sanctified by God’s Spirit.

There is no doubt that God wills to have a living, pure, holy Church, without spot or wrinkle. It is merely an excuse when one speaks of the invisible Church in order to cover over the infirmities of that which we call the congregation among us.

But there is little comfort in this for those who, according to God’s Word, believe that the Church is to be light and salt in the world, the city upon the hill which cannot be hidden. There can be no doubt what the Lord calls his congregation. It is the gathering of all those who through faith in Jesus Christ have become living members of his body, which is the Church.

It is another matter whether here on earth we can bring the Church forth in its true form as a gathering consisting only of believing children of God who have been cleansed and sanctified by the blood of Christ.

We find that this did not succeed even for the Savior himself; for among the little company of disciples there was also a Judas. And the picture that the Apostles’ letters give us of the first congregations also indicates that within them as well there was much that was impure.

But this is nevertheless no reason why we should continue in this country as we have done until now—practically driving our people together into congregations without asking whether there is spiritual life and real Christianity. I do not mention this in order especially to blame particular men or to place the guilt upon any single church body.

No, we all bear our guilt.

Many of us did it in ignorance; we understood no better. In many ways we were zealous in establishing congregations, church bodies, churches, and schools. We gathered people in great masses, large congregations, large church bodies, and we thought that things were going well. But what became of the Lord’s congregation?

We must begin another course of procedure, and let us build from this day onward.

If we are to become true builders, God’s Spirit must be allowed to take up his dwelling in our hearts, and we must not consult flesh and blood as to what I want, or what anyone else wants, but only look to what God’s good and perfect will is. Let us begin in faith.

\subsection{Wednesday Afternoon}

Devotion by Pastor O. Dahle of Aitkin, Minnesota. The times for the business sessions were fixed at 9:30–12:00 in the morning and 2:00–5:00 in the afternoon.

It was resolved that a prayer meeting be held each day from 9:00 to 9:30 in the morning, and that the chairman appoint a man to lead these prayer meetings. It was proposed and resolved that the evening services begin at 8:00.

The chairman was instructed to appoint tellers. It was resolved that the chairman likewise appoint a credentials committee of three members.

The following were appointed tellers: Pastor C. Berntsen, Pastor J. M. Jakobsen, Pastor H. N. Hendrickson, and Pastor J. C. Fossum.

The chairman stated that all voting members of Lutheran congregations, regardless of denominational affiliation, could become voting members of this meeting by signing the following:

“The undersigned hereby declares that he approves the Rules of the Free Church and will work for its purpose as stated in § 2.”

The following were appointed to the credentials committee: Pastor A. J. Løgeland, Prof. M. Pettersen, and Mr. H. Paulson.

It was resolved to hear the reports.

The chairman, Pastor E. P. Harbo of Duluth, first read his report:

At last year’s annual meeting, the Free Church work advanced so far that laws and rules for a Lutheran Free Church were considered and in part adopted. It now falls to this annual meeting to continue the work further. The principles that underlie the work of the Free Church are not new. They are as old as Christianity itself. But among our people it was not until Hans Nielsen Hauge that they began in earnest to make themselves felt. With the Haugean Awakening in Norway there arose a division within the Church, a struggle between the awakened laity and the high-church and, in part, spiritually dead pastors within the State Church.

This struggle has continued here in this country. Those who favor ecclesiastical freedom have here set themselves the goal “to work for the quickening and liberation of all Norwegian Lutheran congregations, so that, according to their calling and ability, they may labor in spiritual independence and freedom for the cause of God’s Kingdom at home and abroad.” This was also established last year as the purpose of the Free Church, and we might readily have expected that all Norwegian Lutheran church bodies in this country would be willing to assent to such a purpose. Nevertheless, we see that they rise up in opposition to it.

The position of the Norwegian Synod in this regard was expressed by Chairman Koren at the Synod’s annual meeting in 1896:

“There is a phenomenon that is well known, though perhaps not so well understood, both in Norway and here in America, and which goes under the name of the western Norwegian Awakening, the sorrowful fruits of which, of various kinds, will remain fresh in the memory and can be recognized year after year by those who have a reasonably exercised discernment and insight into God’s Word. Unfortunately, it has extended much farther than many suppose. From it we will pray God graciously to preserve us. But if we are to have anything capable of standing against this pest, then it must be a true, sound Christian life.”

The United Church consists chiefly of Anti-Missourians who in spirit and disposition are still good Synod people. In this way it can also be explained why the United Church has with such bitterness persecuted those who have labored for awakening and spiritual life within the congregations.

In the United Church it is no longer a question of obedience to conviction and conscience, but first and foremost of obedience to the church body.

The task of the Free Church in this country is therefore plain and clear.

Amid tribulation and opposition, it must labor for freedom of conscience and personal Christian life within the congregations. To this task belongs also the maintenance of a theological seminary whose first aim is to train earnest, living witnesses of the Lord.

Augsburg Seminary has been such a school. In spite of the difficulties, the Seminary has also during the past year been highly blessed by the Lord. Enrollment has been greater than ever before. Eighteen young men graduated from the college department and nine from the theological department. The school, however, has not received the financial support that was needed, which may possibly be due to the unfortunate lawsuit concerning the property and the state of uncertainty into which the school has thereby been brought.

Whatever the outcome of the case may be, it will not have great influence upon the future of the Free Church. Its continued existence does not depend upon earthly property.

The possible loss of the property may, however, make it necessary for this annual meeting to take up for consideration the question of a future home for the theological seminary.

Since the previous annual meeting, Candidate J. Eliasen and Student C. Østerhus have been ordained.

During the past year as well, the work of edification throughout the congregations has been carried on with vigor and has borne good fruit in the conversion and salvation of souls.

Pastor P. Nilsen has again this year traveled about among our people and labored to great blessing. His work should be continued during the coming year as well.

Prof. Sverdrup read the reports of the mission committees, since the mission secretary, Prof. J. H. Blegen, was absent on a journey to Norway. The Committee for Mission to the Heathen had held four meetings during the year. Missionary concern appeared to be increasing strongly, as could be seen, among other things, from the fact that by the close of the year support had been pledged for twenty native teachers.

The committee had resolved to call and send out Candidate J. B. Jerstad as a missionary; likewise, it had already sent out two deaconesses to take charge of the orphanage in place of the late Sister Pernille Pedersen. The sisters at the Deaconess Home had collected about \$300 toward their journey.

Dr. J. Dyrnes and Candidate N. Hatlem had both offered themselves for service as missionaries. Missionary Tou, who is presently at home on a visit, is to travel among the congregations and speak on behalf of the mission cause. During the year, \$5,437.58 had been received into the foreign mission treasury. Approximately \$7,000.00 would be needed for the next year.

The Home Mission Committee had held seven meetings. Its treasury had received \$2,559.95.

At least \$3,000 would be needed for the next year. The committee wished especially to direct the annual meeting’s attention to the missionary need around Fosston, Hawley, and Wahpeton.

The treasurer of the teachers’ salary fund, Prof. J. L. Nydahl, reported that from August 31, 1897, to June 1, 1898, \$3,579.24 had been received for the theological seminary through voluntary contributions.

The Board of Trustees was not yet ready to report; the chairman of the Board, Prof. S. Oftedal, did, however, read the report of the Seminary’s director.

From this it appeared that the theological seminary had this year been blessed by the Lord in a special degree, in spite of the painful position into which the lawsuit over the property had brought it. This unfortunate lawsuit has, however, clearly shown that it is Augsburg’s property, but not Augsburg’s spirit, that the United Church wants.

Even if Augsburg’s property is lost, the work cannot therefore cease. But we must reconcile ourselves to the fact that Augsburg’s work for living pastors and living congregations will meet fierce opposition also in the future, and that it will require faith, love, and perseverance. But then we shall also be permitted to behold the Lord’s help and rejoice in the fruit of our labor, provided that we are found faithful to the last.

The teachers at the school had this year been the same as before; but greater teaching resources and more room were needed. Instruction in English also had to be continued in the theological department.

A class in French had been established this year and led by the director.

The number of students this year had been 194, compared with 187 the year before. Eighteen young men graduated from the college department and nine from the theological department. Of the latter, two had a particular call to become missionaries.

A report from the Free Church Bookstore was read. The business had made good progress. Most of the surplus had been added to the capital fund. An appropriation of \$200.00 had been made to Augsburg Seminary.

The chairman read letters to the annual meeting from Pastors M. Falk Gjertsen and E. Gynild, as well as Prof. Reimestad, all of whom were traveling in Norway. A letter from Pastor O. J. Edwards, Cooperstown, North Dakota, was also read.

It was resolved to proceed to the election of officers. Pastor Harbo was unanimously reelected chairman.

The following were nominated for secretary: Pastor H. Yderstad, Pastor A. Helland, and Pastor H. Urseth. Yderstad received 16 votes, Helland 91, and Urseth 22.

Helland was declared elected. He stated that he could not well accept the election. It was then determined to elect an assistant secretary, and Pastor Yderstad was elected assistant secretary. Yderstad resigned, whereupon the question arose: How, then, are the railroad certificates to be endorsed? After some discussion concerning this, Pastor Yderstad was instructed to investigate the matter. Meanwhile, Pastor Helland asked leave to accept the election as secretary. The resignation was not accepted.

The following were nominated as members of the nominating committee: Prof. J. L. Nydahl, Pastor D. Paulson, Pastor E. Aas, Prof. G. Sverdrup, Mr. Jonas Hetland, Mr. Baard Andersen, and Fingar Fosliet.

Nydahl was elected but resigned. Paulson, Aas, Sverdrup, Hetland, and Andersen were elected. Sverdrup resigned, and Fingar Fosliet was elected in his place.

On the motion of Prof. Sverdrup, it was resolved to take up the “Rules for the Free Church” for consideration on Thursday morning.

In the evening, Pastors A. Halvorson and M. B. Sveen preached in Trinity Church, and Pastors E. Aas and N. J. Thomasberg in St. Luke’s Church.

Continued on page 4. 
